\documentclass[12pt]{article}
\usepackage[utf8]{inputenc}
\usepackage[T1]{fontenc}
\usepackage[margin=1in]{geometry}
\usepackage{setspace}
\usepackage{xcolor}
\usepackage[colorlinks=true, urlcolor=blue!60!black, linkcolor=blue!60!black]{hyperref}
\setlength{\parindent}{0em}
\setlength{\parskip}{0.45em}
\renewcommand{\baselinestretch}{1.0}
\begin{document}
\begin{center}
{\Large \textit{National Republican Senatorial Committee v. Federal Election Commission}, 609 U.S. \_\_\_ (2026)\par}
\vspace{0.3em}
{\normalsize Case Brief --- One-Page Edition\par}
\end{center}

\textbf{Memory Jogger.} Party coordinated-expenditure limits fall as a disproportionate ``fourth prophylaxis''; \textit{Colorado II} overruled (\href{https://www.supremecourt.gov/opinions/25pdf/24-621_h315.pdf}{Slip Opinion}).

\textbf{Facts.} FECA capped party spending in ``coordination'' with federal candidates, as upheld in \textit{Colorado II} (2001). Republican Senate and House committees, JD Vance, and Steve Chabot challenged the caps.

\textbf{Procedural History.} The en banc Sixth Circuit upheld the caps as bound by \textit{Colorado II}. The government declined to defend them; the Supreme Court reversed.

\textbf{Votes.} 6--3: Kavanaugh for the Court; Kagan, with Sotomayor and Jackson, dissented.

\textbf{Holding.} The limits violate the First Amendment. Even under closely drawn scrutiny, a restriction must target \textit{quid pro quo} corruption and be necessary, narrowly tailored, and not disproportionate.

\textbf{Reasoning.} Base limits, earmarking rules, and disclosure already prevent circumvention; \textit{Colorado II} was a ``three-legged stool where all three legs have already been knocked out.'' Kagan answered that joint fundraising lets one donor give \$551,300 for a single candidate's bills.

\textbf{Examples.} \emph{Same-Side:} state caps on party coordination. \emph{Opposite-Side:} Super PAC coordination limits (reserved). \emph{Fence-Sitter:} the cap on direct party cash transfers.

\textbf{Critique.} Some reformers welcome stronger parties; critics say the Court ignored modern fundraising and \textit{stare decisis}.

\textbf{Key Quotation.} ``To sum up: In light of the other meaningful prophylactic measures available to the Government, and given the severe infringement on First Amendment-protected political speech that ensues from limiting a political party's spending in support of its candidates, we conclude that the political-party coordinated-expenditure limits are `disproportionate' and are not `necessary' and `narrowly tailored' for the circumvention interest it seeks to protect.'' --- Justice Kavanaugh, slip op. at 21.

\end{document}
